\documentclass[french,11pt]{book}
\usepackage{teach}
\usepackage[french]{babel}
\usepackage{amsmath,amssymb}
\usepackage{eurosym}
\newcommand{\pourcent}{\%}

\begin{document}

\chapter*{Exercices -- Taux d'evolution, pourcentages et informations chiffrees}

\begin{enumerate}
\item Dans une classe de $32$ eleves, $12$ eleves suivent l'option management. Calculer la proportion correspondante, puis l'exprimer en pourcentage.

\item Un magasin applique une reduction de $18\pourcent$ sur un article coutant $75$ \euro{}. Calculer le montant de la reduction, puis le nouveau prix.

\item Le chiffre d'affaires d'une entreprise passe de $48\,000$ \euro{} a $52\,800$ \euro{}. Calculer le taux d'evolution en pourcentage.

\item Une population diminue de $7\pourcent$. Donner le coefficient multiplicateur associe, puis calculer la nouvelle population si elle etait de $12\,500$ habitants au depart.

\item Un prix subit une hausse de $15\pourcent$ puis une baisse de $15\pourcent$. L'evolution globale est-elle nulle ? Justifier par un calcul.

\item Une valeur est multipliee par $1{,}06$ puis par $0{,}92$. Determiner le taux d'evolution global.

\item Un produit coute $96$ \euro{} apres une augmentation de $20\pourcent$. Retrouver son prix initial.

\item Un salaire augmente de $4\pourcent$ par an pendant trois ans. Le salaire initial est $1\,850$ \euro{}. Calculer le salaire apres trois ans, arrondi au centime.
\end{enumerate}

\end{document}
